% Options for packages loaded elsewhere
\PassOptionsToPackage{unicode}{hyperref}
\PassOptionsToPackage{hyphens}{url}
%
\documentclass[
]{article}
\usepackage{amsmath,amssymb}
\usepackage{iftex}
\ifPDFTeX
  \usepackage[T1]{fontenc}
  \usepackage[utf8]{inputenc}
  \usepackage{textcomp} % provide euro and other symbols
\else % if luatex or xetex
  \usepackage{unicode-math} % this also loads fontspec
  \defaultfontfeatures{Scale=MatchLowercase}
  \defaultfontfeatures[\rmfamily]{Ligatures=TeX,Scale=1}
\fi
\usepackage{lmodern}
\ifPDFTeX\else
  % xetex/luatex font selection
\fi
% Use upquote if available, for straight quotes in verbatim environments
\IfFileExists{upquote.sty}{\usepackage{upquote}}{}
\IfFileExists{microtype.sty}{% use microtype if available
  \usepackage[]{microtype}
  \UseMicrotypeSet[protrusion]{basicmath} % disable protrusion for tt fonts
}{}
\makeatletter
\@ifundefined{KOMAClassName}{% if non-KOMA class
  \IfFileExists{parskip.sty}{%
    \usepackage{parskip}
  }{% else
    \setlength{\parindent}{0pt}
    \setlength{\parskip}{6pt plus 2pt minus 1pt}}
}{% if KOMA class
  \KOMAoptions{parskip=half}}
\makeatother
\usepackage{xcolor}
\usepackage{color}
\usepackage{fancyvrb}
\newcommand{\VerbBar}{|}
\newcommand{\VERB}{\Verb[commandchars=\\\{\}]}
\DefineVerbatimEnvironment{Highlighting}{Verbatim}{commandchars=\\\{\}}
% Add ',fontsize=\small' for more characters per line
\newenvironment{Shaded}{}{}
\newcommand{\AlertTok}[1]{\textcolor[rgb]{1.00,0.00,0.00}{\textbf{#1}}}
\newcommand{\AnnotationTok}[1]{\textcolor[rgb]{0.38,0.63,0.69}{\textbf{\textit{#1}}}}
\newcommand{\AttributeTok}[1]{\textcolor[rgb]{0.49,0.56,0.16}{#1}}
\newcommand{\BaseNTok}[1]{\textcolor[rgb]{0.25,0.63,0.44}{#1}}
\newcommand{\BuiltInTok}[1]{\textcolor[rgb]{0.00,0.50,0.00}{#1}}
\newcommand{\CharTok}[1]{\textcolor[rgb]{0.25,0.44,0.63}{#1}}
\newcommand{\CommentTok}[1]{\textcolor[rgb]{0.38,0.63,0.69}{\textit{#1}}}
\newcommand{\CommentVarTok}[1]{\textcolor[rgb]{0.38,0.63,0.69}{\textbf{\textit{#1}}}}
\newcommand{\ConstantTok}[1]{\textcolor[rgb]{0.53,0.00,0.00}{#1}}
\newcommand{\ControlFlowTok}[1]{\textcolor[rgb]{0.00,0.44,0.13}{\textbf{#1}}}
\newcommand{\DataTypeTok}[1]{\textcolor[rgb]{0.56,0.13,0.00}{#1}}
\newcommand{\DecValTok}[1]{\textcolor[rgb]{0.25,0.63,0.44}{#1}}
\newcommand{\DocumentationTok}[1]{\textcolor[rgb]{0.73,0.13,0.13}{\textit{#1}}}
\newcommand{\ErrorTok}[1]{\textcolor[rgb]{1.00,0.00,0.00}{\textbf{#1}}}
\newcommand{\ExtensionTok}[1]{#1}
\newcommand{\FloatTok}[1]{\textcolor[rgb]{0.25,0.63,0.44}{#1}}
\newcommand{\FunctionTok}[1]{\textcolor[rgb]{0.02,0.16,0.49}{#1}}
\newcommand{\ImportTok}[1]{\textcolor[rgb]{0.00,0.50,0.00}{\textbf{#1}}}
\newcommand{\InformationTok}[1]{\textcolor[rgb]{0.38,0.63,0.69}{\textbf{\textit{#1}}}}
\newcommand{\KeywordTok}[1]{\textcolor[rgb]{0.00,0.44,0.13}{\textbf{#1}}}
\newcommand{\NormalTok}[1]{#1}
\newcommand{\OperatorTok}[1]{\textcolor[rgb]{0.40,0.40,0.40}{#1}}
\newcommand{\OtherTok}[1]{\textcolor[rgb]{0.00,0.44,0.13}{#1}}
\newcommand{\PreprocessorTok}[1]{\textcolor[rgb]{0.74,0.48,0.00}{#1}}
\newcommand{\RegionMarkerTok}[1]{#1}
\newcommand{\SpecialCharTok}[1]{\textcolor[rgb]{0.25,0.44,0.63}{#1}}
\newcommand{\SpecialStringTok}[1]{\textcolor[rgb]{0.73,0.40,0.53}{#1}}
\newcommand{\StringTok}[1]{\textcolor[rgb]{0.25,0.44,0.63}{#1}}
\newcommand{\VariableTok}[1]{\textcolor[rgb]{0.10,0.09,0.49}{#1}}
\newcommand{\VerbatimStringTok}[1]{\textcolor[rgb]{0.25,0.44,0.63}{#1}}
\newcommand{\WarningTok}[1]{\textcolor[rgb]{0.38,0.63,0.69}{\textbf{\textit{#1}}}}
\usepackage{longtable,booktabs,array}
\usepackage{calc} % for calculating minipage widths
% Correct order of tables after \paragraph or \subparagraph
\usepackage{etoolbox}
\makeatletter
\patchcmd\longtable{\par}{\if@noskipsec\mbox{}\fi\par}{}{}
\makeatother
% Allow footnotes in longtable head/foot
\IfFileExists{footnotehyper.sty}{\usepackage{footnotehyper}}{\usepackage{footnote}}
\makesavenoteenv{longtable}
\setlength{\emergencystretch}{3em} % prevent overfull lines
\providecommand{\tightlist}{%
  \setlength{\itemsep}{0pt}\setlength{\parskip}{0pt}}
\setcounter{secnumdepth}{-\maxdimen} % remove section numbering
\ifLuaTeX
  \usepackage{selnolig}  % disable illegal ligatures
\fi
\usepackage[]{natbib}
\bibliographystyle{plainnat}
\IfFileExists{bookmark.sty}{\usepackage{bookmark}}{\usepackage{hyperref}}
\IfFileExists{xurl.sty}{\usepackage{xurl}}{} % add URL line breaks if available
\urlstyle{same}
\hypersetup{
  hidelinks,
  pdfcreator={LaTeX via pandoc}}

\author{}
\date{}

\begin{document}

\hypertarget{when-does-dynamic-weight-decay-help-a-unified-framework-analysis-under-adamw}{%
\section{When Does Dynamic Weight Decay Help? A Unified Framework
Analysis Under
AdamW}\label{when-does-dynamic-weight-decay-help-a-unified-framework-analysis-under-adamw}}

\begin{center}\rule{0.5\linewidth}{0.5pt}\end{center}

\hypertarget{abstract}{%
\subsection{Abstract}\label{abstract}}

Weight decay is universally applied in deep learning, yet the scheduling
and modulation of its strength throughout training remains poorly
understood. Recent proposals---including alignment-aware decay (Cautious
Weight Decay), gradient-norm scheduling (Scheduled Weight Decay), cosine
annealing, and spatial modulation---each claim improvements over the
constant baseline, but are evaluated under incompatible experimental
conditions, making direct comparison impossible. We introduce the
\textbf{Phi Modulator Framework}, a unified mathematical abstraction
\(\varphi(t, \boldsymbol{\theta}, \mathbf{g}) : \mathbb{Z}_{\geq 0} \times \mathbb{R}^d \times \mathbb{R}^d \to \mathbb{R}^d_{\geq 0}\)
that subsumes all major dynamic weight decay strategies as special cases
along four modulation axes: temporal, directional, spatial, and
target-norm. Alongside the framework we propose three standardized
diagnostic metrics---the Budget Equivalence Metric (BEM), Coupling
Stability Index (CSI), and Alignment Informativeness Score
(AIS)---providing the first quantitative tools for characterizing how
weight decay methods differ in effective budget, optimization coupling,
and alignment informativeness. In a systematic evaluation of seven
weight decay strategies on CIFAR-10 and CIFAR-100 with ResNet-20 under
AdamW (42 experiments, three seeds per configuration), we find that
\textbf{all dynamic weight decay variants are statistically
indistinguishable from constant weight decay under AdamW (\(p > 0.05\)),
including the degenerate case of no weight decay}. A ten-fold variation
in effective weight decay budget (BEM range \(0.0\)--\(1.0\)) produces
less than \(0.5\%\) accuracy variation. We formalize this as the
\emph{Phi Invariance Conjecture} scoped to CIFAR-scale BatchNorm
networks under AdamW with moderate \(\lambda\): AdamW's per-parameter
adaptive scaling subsumes the effect of any Phi modulator. Critically,
49 SGD control experiments on the same configurations confirm the
conjecture is optimizer-specific: under SGD, removing weight decay
entirely causes a statistically significant \(-0.91\%\) accuracy drop
(\(p < 0.001\)), and gradient-norm scheduling (SWD) degrades performance
by \(-0.51\%\) (\(p = 0.013\)). This SGD contrast validates that the
AdamW invariance is a genuine property of adaptive optimization, not a
statistical artifact. The framework and diagnostic metrics provide the
first standardized infrastructure for weight decay research.

\begin{center}\rule{0.5\linewidth}{0.5pt}\end{center}

\hypertarget{introduction}{%
\subsection{1. Introduction}\label{introduction}}

\hypertarget{motivation}{%
\subsubsection{1.1 Motivation}\label{motivation}}

Weight decay is among the most universally applied techniques in deep
learning optimization. Virtually every modern training recipe---from
small-scale CIFAR classifiers to billion-parameter language
models---includes a weight decay coefficient as a core hyperparameter.
Yet despite its ubiquity, the community lacks a principled framework for
choosing \emph{how} to apply weight decay over the course of training.

The classical understanding treats weight decay as explicit L2
regularization that shrinks weights toward zero, discouraging model
complexity (Krogh \& Hertz, 1991). This view has been progressively
undermined by modern findings. Loshchilov \& Hutter (2019) demonstrated
that L2 regularization and weight decay are not equivalent in adaptive
optimizers, leading to the now-standard AdamW formulation that
\emph{decouples} weight decay from gradient scaling. More recently,
D'Angelo et al.~(2024) showed that weight decay is never useful as
explicit regularization in modern deep learning; instead, it acts as a
\emph{training dynamics modifier}---stabilizing loss trajectories under
SGD and controlling the bias-variance tradeoff in large language model
training.

This re-understanding has sparked a surge of methods proposing to
dynamically modulate weight decay strength during training. Xie et
al.~(2023) introduced Scheduled Weight Decay (SWD), which adjusts weight
decay based on gradient norms. Chen et al.~(2026a; CWD) proposed
Cautious Weight Decay, a sign-alignment mask that applies decay only
when weight and update directions agree. Loshchilov (2023) generalized
decoupled weight decay to target-norm control (AdamWN), while He et
al.~(2025) introduced AlphaDecay, a module-wise strategy guided by
spectral density analysis. Ferbach et al.~(2026) proposed ADANA,
applying logarithmic-time schedules to both weight decay and momentum.
Chen et al.~(2026b; AdamO) identified a ``Radial Tug-of-War'' conflict
between weight decay and gradient updates.

A critical problem pervades this rapidly growing literature:
\textbf{each method is evaluated in isolation}, using different
architectures, datasets, optimizers, hyperparameter selection protocols,
and evaluation metrics. No two papers share the same experimental
conditions, making direct comparison impossible.

This raises a fundamental question: \textbf{does dynamic weight decay
actually help, and if so, when and why?}

\hypertarget{research-gap}{%
\subsubsection{1.2 Research Gap}\label{research-gap}}

Four critical gaps prevent answering this question:

\textbf{Gap 1: No unified mathematical framework.} The four major
families of dynamic weight decay---temporal scheduling, directional
modulation, spatial modulation, and target-norm control---each operate
with independent mathematical formulations. No existing work reveals
whether they are fundamentally different or special cases of a single
principle.

\textbf{Gap 2: No standardized evaluation metrics.} Each paper reports
different metrics under different conditions. There is no standard way
to quantify how much effective weight decay budget a method uses, how
stably it couples with the optimizer, or whether its alignment signal
carries useful information.

\textbf{Gap 3: No controlled systematic comparison.} No prior work has
evaluated multiple dynamic weight decay methods within a single codebase
under identical hyperparameters with proper statistical testing.

\textbf{Gap 4: No theory for when dynamic weight decay matters.}
D'Angelo et al.~(2024) showed weight decay acts as a dynamics modifier,
implying scheduling should matter. Yet no theoretical framework predicts
\emph{when} a particular strategy outperforms constant weight decay. Xie
\& Li (2024) showed AdamW implicitly performs \(\ell_\infty\)-norm
constrained optimization, hinting at a potential absorption mechanism,
but this has not been connected to weight decay scheduling.

\hypertarget{contributions}{%
\subsubsection{1.3 Contributions}\label{contributions}}

\begin{enumerate}
\def\labelenumi{\arabic{enumi}.}
\item
  \textbf{The Phi Modulator Framework.} A unified mathematical interface
  \(\varphi(t, \boldsymbol{\theta}, \mathbf{g})\) that expresses the
  weight decay update as
  \(\boldsymbol{\theta}_{t+1} = \boldsymbol{\theta}_t - \eta_t \hat{\mathbf{m}}_t / (\sqrt{\hat{\mathbf{v}}_t} + \epsilon) - \lambda \cdot \varphi(t, \boldsymbol{\theta}_t, \mathbf{g}_t) \odot \boldsymbol{\theta}_t\).
  CWD, SWD, cosine schedules, AdamWN, AlphaDecay, and all compositions
  are recovered as special cases along four modulation axes.
\item
  \textbf{Three diagnostic metrics.} BEM, CSI, and AIS---the first
  standardized tools for quantifying how weight decay methods differ in
  effective budget, coupling stability, and alignment informativeness,
  even when they produce identical final accuracy.
\item
  \textbf{Systematic benchmark with SGD negative control.} 42 controlled
  AdamW experiments (7 methods \(\times\) 3 seeds \(\times\) 2 datasets)
  plus 49 SGD control experiments on the same configurations. The SGD
  experiments serve as a negative control: if the null result were a
  CIFAR-scale power artifact, SGD would also show invariance. It does
  not.
\item
  \textbf{The Phi Invariance Conjecture (scoped).} Under AdamW with
  moderate \(\lambda\) on CIFAR-scale BatchNorm networks, all dynamic
  weight decay variants are statistically equivalent to constant weight
  decay. The SGD contrast validates that this invariance is a genuine
  property of AdamW's adaptive mechanism: weight decay magnitude matters
  significantly under SGD (no\_wd: \(-0.91\%\), \(p < 0.001\); SWD:
  \(-0.51\%\), \(p = 0.013\)), but not under AdamW.
\end{enumerate}

\hypertarget{paper-roadmap}{%
\subsubsection{1.4 Paper Roadmap}\label{paper-roadmap}}

Section 2 surveys related work. Section 3 introduces the Phi Modulator
Framework and diagnostic metrics. Section 4 describes the experimental
setup (AdamW and SGD). Section 5 presents AdamW results and diagnostic
analysis. Section 6 presents the Phi Invariance Conjecture with SGD
validation and boundary conditions. Section 7 concludes.

\begin{center}\rule{0.5\linewidth}{0.5pt}\end{center}

\hypertarget{related-work}{%
\subsection{2. Related Work}\label{related-work}}

\hypertarget{weight-decay-as-a-dynamics-modifier}{%
\subsubsection{2.1 Weight Decay as a Dynamics
Modifier}\label{weight-decay-as-a-dynamics-modifier}}

The classical view treats weight decay as L2 regularization (Krogh \&
Hertz, 1991). Loshchilov \& Hutter (2019) demonstrated that in adaptive
optimizers, L2 regularization and decoupled weight decay produce
fundamentally different behaviors. Their proposed AdamW has since become
the default optimizer for modern deep learning.

D'Angelo et al.~(2024) showed weight decay is \emph{never} useful as
explicit regularization in modern settings; it serves as a training
dynamics modifier. Kosson et al.~(2023) showed weight decay under AdamW
induces a \emph{rotational equilibrium} that balances weight vector
rotation across layers. If this equilibrium is achieved robustly
regardless of the specific Phi modulation form, then the scheduling
strategy becomes irrelevant---a hypothesis our experiments test. Xie \&
Li (2024) proved AdamW implicitly performs \(\ell_\infty\)-norm
constrained optimization, connecting decoupled weight decay to the
Frank-Wolfe algorithm.

\hypertarget{dynamic-weight-decay-methods}{%
\subsubsection{2.2 Dynamic Weight Decay
Methods}\label{dynamic-weight-decay-methods}}

\textbf{Temporal scheduling.} Xie et al.~(2023) introduced Scheduled
Weight Decay (SWD/AdamS), adjusting weight decay based on gradient
norms. Ferbach et al.~(2026) proposed ADANA, applying logarithmic-time
schedules to weight decay alongside momentum coefficients.

\textbf{Directional modulation.} Chen et al.~(2026a) proposed CWD,
applying a binary sign-alignment mask. The CWD paper reports
improvements primarily with Lion, Muon, and language model
pre-training---not under AdamW at CIFAR scale. Chen et al.~(2026b)
proposed AdamO to decouple radial and tangential dynamics. Tian et
al.~(2024) introduced SPD for fine-tuning.

\textbf{Spatial modulation.} He et al.~(2025) proposed AlphaDecay,
assigning module-wise decay rates guided by heavy-tailed
self-regularization spectral density analysis, demonstrating gains at
LLM scales from 60M to 1B parameters.

\textbf{Target-norm control.} Loshchilov (2023) generalized decoupled
weight decay to Weight Norm Control (AdamWN), driving parameters toward
an arbitrary target norm \(\tau\). Wang \& Aitchison (2024) showed
optimal weight decay scales as an EMA timescale constant, suggesting a
well-calibrated constant weight decay may already capture the available
benefit.

\textbf{Structural effects.} Galanti et al.~(2022) demonstrated that SGD
with weight decay induces a low-rank bias. Kobayashi et al.~(2024)
showed L2 regularization on multiplicative parameters is equivalent to
nuclear norm regularization.

\hypertarget{evaluation-fragmentation}{%
\subsubsection{2.3 Evaluation
Fragmentation}\label{evaluation-fragmentation}}

CWD is evaluated with Lion and Muon on language model pre-training; SWD
targets the SGD-Adam generalization gap on CIFAR and ImageNet;
AlphaDecay operates at LLM scales. No two papers share the same
experimental conditions. Fernandez-Hernandez et al.~(2025) proposed the
Overfitting-Underfitting Indicator (OUI) as a per-method diagnostic.
D'Angelo et al.~(2024) provided shared experimental infrastructure but
did not compare dynamic scheduling strategies. This fragmentation
motivates our standardized metrics and systematic benchmark.

\begin{center}\rule{0.5\linewidth}{0.5pt}\end{center}

\hypertarget{the-phi-modulator-framework}{%
\subsection{3. The Phi Modulator
Framework}\label{the-phi-modulator-framework}}

\hypertarget{formal-definition}{%
\subsubsection{3.1 Formal Definition}\label{formal-definition}}

Consider a neural network with parameters
\(\boldsymbol{\theta} \in \mathbb{R}^d\) trained with AdamW. The
standard AdamW update rule at step \(t\) is: \[
\boldsymbol{\theta}_{t+1} = \boldsymbol{\theta}_t - \eta_t \frac{\hat{\mathbf{m}}_t}{\sqrt{\hat{\mathbf{v}}_t} + \epsilon} - \lambda \boldsymbol{\theta}_t
\] where \(\hat{\mathbf{m}}_t\) and \(\hat{\mathbf{v}}_t\) are the
bias-corrected first and second moment estimates, \(\eta_t\) is the
learning rate, \(\epsilon\) is a stability constant, and \(\lambda\) is
the weight decay coefficient. We generalize this by introducing the
\textbf{Phi modulator} \(\varphi\): \[
\boldsymbol{\theta}_{t+1} = \boldsymbol{\theta}_t - \eta_t \frac{\hat{\mathbf{m}}_t}{\sqrt{\hat{\mathbf{v}}_t} + \epsilon} - \lambda \cdot \varphi(t, \boldsymbol{\theta}_t, \mathbf{g}_t) \odot \boldsymbol{\theta}_t
\] where
\(\varphi : \mathbb{Z}_{\geq 0} \times \mathbb{R}^d \times \mathbb{R}^d \to \mathbb{R}^d_{\geq 0}\)
is a per-parameter, non-negative modulation function and \(\odot\)
denotes element-wise multiplication. We define
\(\mathbf{u}_t = \hat{\mathbf{m}}_t / (\sqrt{\hat{\mathbf{v}}_t} + \epsilon)\)
as the \textbf{optimizer update direction}. Modulators that condition on
alignment (e.g., CWD) use \(\mathbf{u}_t\) rather than the raw gradient
\(\mathbf{g}_t\).

The Phi modulator satisfies three key properties: - \textbf{Positivity:}
\(\varphi(t, \boldsymbol{\theta}, \mathbf{g}) \geq 0\) component-wise. -
\textbf{Measurability:} \(\varphi\) can depend on training step,
parameters, gradients, and optimizer state. - \textbf{Normalization
reference:} \(\mathbb{E}[\varphi] = 1\) is the reference target
(budget-equivalent to constant weight decay). Deviations are quantified
by BEM (Section 3.4).

The framework admits a clean programmatic interface:

\begin{Shaded}
\begin{Highlighting}[]
\KeywordTok{class}\NormalTok{ WDModulator(ABC):}
    \KeywordTok{def}\NormalTok{ compute\_phi(}\VariableTok{self}\NormalTok{, w: Tensor, u: Tensor, t: }\BuiltInTok{int}\NormalTok{) }\OperatorTok{{-}\textgreater{}}\NormalTok{ Tensor:}
        \CommentTok{"""Return per{-}parameter modulation phi in [0, inf).}
\CommentTok{        Args:}
\CommentTok{            w: parameter tensor (theta\_t)}
\CommentTok{            u: optimizer update direction (m\_hat / (sqrt(v\_hat) + eps))}
\CommentTok{            t: training step}
\CommentTok{        """}
\NormalTok{        ...}
\end{Highlighting}
\end{Shaded}

Every weight decay strategy in our study is implemented as a subclass of
this interface, ensuring identical integration with the base optimizer.

\hypertarget{special-cases-recovering-existing-methods}{%
\subsubsection{3.2 Special Cases: Recovering Existing
Methods}\label{special-cases-recovering-existing-methods}}

Table 1 organizes the major dynamic weight decay methods along four
modulation axes as special cases of \(\varphi\).

\textbf{Table 1: Method catalog.} Each row specifies the closed-form
\(\varphi\) expression and the modulation axis.

\begin{longtable}[]{@{}
  >{\raggedright\arraybackslash}p{(\columnwidth - 4\tabcolsep) * \real{0.1111}}
  >{\raggedright\arraybackslash}p{(\columnwidth - 4\tabcolsep) * \real{0.6528}}
  >{\raggedright\arraybackslash}p{(\columnwidth - 4\tabcolsep) * \real{0.2361}}@{}}
\toprule\noalign{}
\begin{minipage}[b]{\linewidth}\raggedright
Method
\end{minipage} & \begin{minipage}[b]{\linewidth}\raggedright
\(\varphi(t, \boldsymbol{\theta}, \mathbf{g})\)
\end{minipage} & \begin{minipage}[b]{\linewidth}\raggedright
Modulation Axis
\end{minipage} \\
\midrule\noalign{}
\endhead
\bottomrule\noalign{}
\endlastfoot
AdamW (constant) & \(\mathbf{1}\) & Baseline (\(\varphi \equiv 1\)) \\
CWD (hard) &
\(\mathbb{1}[\mathrm{sign}(\boldsymbol{\theta}) = \mathrm{sign}(\mathbf{u}_t)]\)
& Directional \\
CWD (soft, \(\beta\)) &
\(\sigma(\beta \cdot \boldsymbol{\theta} \odot \mathbf{u}_t)\) &
Directional \\
SWD / AdamS & \(h(\|\mathbf{g}_t\|) \cdot \mathbf{1}\),
\(h(x) = x / (\bar{x} + \epsilon_h)\) & Temporal-gradient \\
Cosine WD & \(\tfrac{1}{2}(1 + \cos(\pi t / T)) \cdot \mathbf{1}\) &
Temporal \\
AdamWN &
\(\max(0, 1 - \tau / \|\boldsymbol{\theta}\|) \cdot \mathbf{1}\) &
Target-norm \\
AlphaDecay & \(\alpha_l \cdot \mathbf{1}_l\) (per layer \(l\)) &
Spatial \\
No-WD & \(\mathbf{0}\) & Ablation \\
Random mask & \(\mathrm{Bernoulli}(p) \cdot \mathbf{1}\) & Control \\
\end{longtable}

Here
\(\mathbf{u}_t = \hat{\mathbf{m}}_t / (\sqrt{\hat{\mathbf{v}}_t} + \epsilon)\),
\(\sigma(\cdot)\) is the sigmoid function,
\(h(x) = x / (\bar{x} + \epsilon_h)\) is SWD's gradient-norm sensitivity
function (where \(\bar{x}\) is the running mean of
\(\|\mathbf{g}_t\|\)), \(T\) is the total number of training steps,
\(\tau\) is AdamWN's target norm, and \(\alpha_l\) is AlphaDecay's
per-layer decay coefficient.

\textbf{Proposition 1} (Composition). \emph{For any two valid Phi
modulators \(\varphi_1\) and \(\varphi_2\), their element-wise product
\(\varphi_{\mathrm{comp}} = \varphi_1 \odot \varphi_2\) is also a valid
Phi modulator.} This follows directly from the positivity property.
Composition formalizes strategies such as CWD+Cosine and CWD+AdamWN as
principled combinations.

\hypertarget{budget-equivalence-normalization}{%
\subsubsection{3.3 Budget Equivalence
Normalization}\label{budget-equivalence-normalization}}

\textbf{Definition 1} (Budget Equivalence). \emph{Two weight decay
strategies with modulators \(\varphi_1\) and \(\varphi_2\) are
budget-equivalent if:} \[
\sum_{t=0}^{T} \lambda \cdot \mathbb{E}_{\boldsymbol{\theta}}[\varphi_1(t, \boldsymbol{\theta}_t, \mathbf{g}_t)] = \sum_{t=0}^{T} \lambda \cdot \mathbb{E}_{\boldsymbol{\theta}}[\varphi_2(t, \boldsymbol{\theta}_t, \mathbf{g}_t)]
\]

Budget equivalence requires matching the time-averaged effective decay
\(\bar{\lambda}_{\mathrm{eff}} = \frac{1}{T}\sum_{t=0}^{T}\lambda_{\mathrm{eff}}(t)\)
before attributing accuracy differences to the scheduling strategy.

\hypertarget{diagnostic-metrics}{%
\subsubsection{3.4 Diagnostic Metrics}\label{diagnostic-metrics}}

\textbf{Budget Equivalence Metric (BEM).} The BEM quantifies deviation
from the constant baseline's effective budget: \[
\mathrm{BEM}(\text{method}) = \frac{|\bar{\lambda}_{\mathrm{eff}}^{\text{method}} - \bar{\lambda}_{\mathrm{eff}}^{\text{constant}}|}{\bar{\lambda}_{\mathrm{eff}}^{\text{constant}}}
\] \(\mathrm{BEM} = 0\) indicates identical budget to constant;
\(\mathrm{BEM} = 1\) indicates zero effective weight decay. BEM enables
disentangling \emph{how much} weight decay is applied from \emph{how} it
is distributed.

\textbf{Coupling Stability Index (CSI).} The CSI measures the stability
of the coupling between weight decay and optimizer adaptation dynamics:
\[
\mathrm{CSI} = w_1 \cdot \widetilde{\mathrm{CV}}(\|\boldsymbol{\theta}\|_{\text{trajectory}}) + w_2 \cdot \widetilde{\log \kappa}(\mathbf{H}_{\text{final}}) + w_3 \cdot \widetilde{\mathrm{CV}}(\eta_{\mathrm{eff}, \text{layers}})
\] where \(\widetilde{(\cdot)}\) denotes normalization relative to the
constant-baseline value, \(\mathrm{CV}(\cdot)\) is the coefficient of
variation, \(\kappa(\mathbf{H}_{\text{final}})\) is the spectral
condition number (approximated via power iteration), and
\(\eta_{\mathrm{eff}} = \eta / (1 + \lambda \|\boldsymbol{\theta}_l\|)\).
Component weights: \(w_1 = 0.4\), \(w_2 = 0.3\), \(w_3 = 0.3\); a
sensitivity analysis (Appendix C) confirms CSI rankings are robust to
\(w_1 \in [0.3, 0.5]\). Higher CSI indicates more unstable coupling.

\textbf{Alignment Informativeness Score (AIS).} The AIS measures whether
the geometric alignment between weights and gradients carries predictive
signal for training progress: \[
\mathrm{AIS} = \frac{1}{L}\sum_{l=1}^{L} |\rho_S|\!\left(\cos(\boldsymbol{\theta}^{(l)}_i, \mathbf{g}^{(l)}_i),\; \Delta\mathcal{L}_i\right)
\] where \(\rho_S\) is the Spearman rank correlation and
\(\Delta\mathcal{L}_i = \mathcal{L}(\boldsymbol{\theta}_{i}) - \mathcal{L}(\boldsymbol{\theta}_{i+1})\)
is the per-step loss reduction. \(\mathrm{AIS} \in [0, 1]\). The
threshold \(\mathrm{AIS} > 0.2\) for ``informative alignment'' is
calibrated against the random\_mask control in our experiments: the
random\_mask baseline (which applies Bernoulli weight decay independent
of alignment) achieves \(\mathrm{AIS} = 0.359 \pm 0.03\), confirming
that values in this range reflect genuine geometric structure in the
loss landscape rather than noise.

\begin{center}\rule{0.5\linewidth}{0.5pt}\end{center}

\hypertarget{experimental-setup}{%
\subsection{4. Experimental Setup}\label{experimental-setup}}

\hypertarget{implementation}{%
\subsubsection{4.1 Implementation}\label{implementation}}

All experiments use a unified \texttt{UnifiedAdamW} (or
\texttt{UnifiedSGD}) optimizer class with a pluggable Phi modulator
interface in PyTorch. Each weight decay strategy is a subclass of
\texttt{WDModulator} (Section 3.1), ensuring all methods share identical
optimizer internals and differ only in \(\varphi\). The codebase extends
the \texttt{why-weight-decay} infrastructure (D'Angelo et al., 2024).

Seven methods spanning the four modulation axes: \textbf{constant}
(baseline, \(\varphi \equiv 1\)), \textbf{cwd\_hard} (Cautious Weight
Decay), \textbf{swd} (Scheduled Weight Decay), \textbf{cosine\_schedule}
(cosine-annealed weight decay), \textbf{random\_mask} (Bernoulli mask,
\(p = 0.5\)), \textbf{half\_lambda} (constant weight decay at
\(\lambda/2\)), and \textbf{no\_wd} (\(\varphi \equiv 0\)).

\hypertarget{training-configuration}{%
\subsubsection{4.2 Training
Configuration}\label{training-configuration}}

\textbf{Datasets.} CIFAR-10 and CIFAR-100 (Krizhevsky, 2009), standard
splits. Augmentation: random cropping (4-pixel padding), random
horizontal flip, per-channel normalization.

\textbf{Architecture.} ResNet-20 (\$\sim\$270K parameters) with batch
normalization (He et al., 2016).

\textbf{AdamW optimizer.} \(\eta = 10^{-3}\) with cosine annealing to
zero (no warmup), \(\beta_1 = 0.9\), \(\beta_2 = 0.999\),
\(\epsilon = 10^{-8}\).

\textbf{SGD optimizer (negative control).} SGD with momentum 0.9,
\(\eta = 0.1\) with cosine annealing, same seven methods and base weight
decay \(\lambda = 5 \times 10^{-4}\). This design is deliberate: if the
Phi Invariance Conjecture is specific to AdamW's adaptive mechanism, SGD
should not exhibit it.

\textbf{Weight decay.} \(\lambda = 5 \times 10^{-4}\) for all methods.

\textbf{Training duration.} 200 epochs, batch size 128.

\textbf{Seeds.} Three independent runs per method--dataset--optimizer
configuration using seeds \(\{42, 123, 456\}\). Total: 42 AdamW
experiments + 49 SGD experiments (3 seeds for 5 methods on CIFAR-100
where complete; 3 seeds for all 7 methods on CIFAR-10).

\hypertarget{hyperparameter-fairness-protocol}{%
\subsubsection{4.3 Hyperparameter Fairness
Protocol}\label{hyperparameter-fairness-protocol}}

\textbf{All methods use identical base hyperparameters} with no
per-method grid search. This ensures observed differences measure the
Phi modulation effect, not hyperparameter luck. Each method may operate
at a non-optimal configuration, but this is the price of a fair
comparison.

\hypertarget{diagnostic-logging}{%
\subsubsection{4.4 Diagnostic Logging}\label{diagnostic-logging}}

We log per-epoch: test accuracy, training loss, per-layer weight norms,
CSI, AIS, and BEM. Full diagnostic panels for all 42 AdamW runs are
provided in Appendix B.

\hypertarget{statistical-power}{%
\subsubsection{4.5 Statistical Power}\label{statistical-power}}

With \(N = 3\) seeds per configuration, our paired \(t\)-tests have 2
degrees of freedom. At 80\% power and \(\alpha = 0.05\) (two-tailed),
the minimum detectable effect size is approximately 0.7\%, given
observed within-method standard deviations of \$\sim\$0.3\%.

For a null-result paper, failure to reject is insufficient. We apply Two
One-Sided Tests (TOST) equivalence testing for the AdamW results
(Section 5.1) and use the SGD negative control as a design-level
validation that does not rely on equivalence testing power (Section
6.1).

\begin{center}\rule{0.5\linewidth}{0.5pt}\end{center}

\hypertarget{results-and-analysis-adamw}{%
\subsection{5. Results and Analysis
(AdamW)}\label{results-and-analysis-adamw}}

\hypertarget{main-accuracy-comparison}{%
\subsubsection{5.1 Main Accuracy
Comparison}\label{main-accuracy-comparison}}

Figure 2 and Table 2 present the AdamW results. On CIFAR-10, the seven
methods span a total accuracy range of 0.25 percentage points: from
89.88\% (SWD) to 90.13\% (constant). On CIFAR-100, the range is 0.76
percentage points: from 62.66\% (no\_wd) to 63.42\% (cosine\_schedule).
Neither extreme is statistically significant.

\textbf{Table 2: Main accuracy results (AdamW).} Mean \(\pm\) standard
deviation over 3 seeds. Best result per dataset in \textbf{bold}.

\begin{longtable}[]{@{}lcc@{}}
\toprule\noalign{}
Method & CIFAR-10 Acc. (\%) & CIFAR-100 Acc. (\%) \\
\midrule\noalign{}
\endhead
\bottomrule\noalign{}
\endlastfoot
constant & \textbf{90.13} \(\pm\) 0.31 & 63.15 \(\pm\) 0.30 \\
cosine\_schedule & 90.12 \(\pm\) \textbf{0.07} & \textbf{63.42} \(\pm\)
0.42 \\
random\_mask & 90.12 \(\pm\) 0.30 & 62.87 \(\pm\) 0.38 \\
half\_lambda & 90.09 \(\pm\) 0.29 & 62.91 \(\pm\) 0.47 \\
no\_wd & 90.08 \(\pm\) 0.32 & 62.66 \(\pm\) 0.38 \\
cwd\_hard & 90.06 \(\pm\) 0.24 & 62.84 \(\pm\) 0.30 \\
swd & 89.88 \(\pm\) 0.25 & 63.06 \(\pm\) 0.29 \\
\end{longtable}

cosine\_schedule achieves the lowest variance on CIFAR-10
(\(\sigma = 0.07\%\), vs.~\(\sigma \approx 0.25\)--\(0.32\%\) for all
other methods), suggesting a potential stability benefit without an
accuracy advantage (Section 6.4).

Table 3 reports paired \(t\)-tests against the constant baseline. All
\(p\)-values exceed 0.05 (range: \(p = 0.090\)--\(0.950\)). After
Bonferroni correction, no method survives. All Cohen's \(d\) effect
sizes are below 0.3.

\textbf{Table 3: Statistical tests vs.~constant baseline (AdamW).}
\(\Delta\) = method accuracy minus constant accuracy.

\begin{longtable}[]{@{}
  >{\raggedright\arraybackslash}p{(\columnwidth - 8\tabcolsep) * \real{0.1159}}
  >{\centering\arraybackslash}p{(\columnwidth - 8\tabcolsep) * \real{0.2754}}
  >{\centering\arraybackslash}p{(\columnwidth - 8\tabcolsep) * \real{0.1594}}
  >{\centering\arraybackslash}p{(\columnwidth - 8\tabcolsep) * \real{0.2899}}
  >{\centering\arraybackslash}p{(\columnwidth - 8\tabcolsep) * \real{0.1594}}@{}}
\toprule\noalign{}
\begin{minipage}[b]{\linewidth}\raggedright
Method
\end{minipage} & \begin{minipage}[b]{\linewidth}\centering
CIFAR-10 \(\Delta\)
\end{minipage} & \begin{minipage}[b]{\linewidth}\centering
\(p\)-value
\end{minipage} & \begin{minipage}[b]{\linewidth}\centering
CIFAR-100 \(\Delta\)
\end{minipage} & \begin{minipage}[b]{\linewidth}\centering
\(p\)-value
\end{minipage} \\
\midrule\noalign{}
\endhead
\bottomrule\noalign{}
\endlastfoot
cwd\_hard & \(-0.07\%\) & 0.832 & \(-0.31\%\) & 0.326 \\
swd & \(-0.25\%\) & 0.513 & \(-0.10\%\) & 0.801 \\
cosine\_schedule & \(-0.01\%\) & 0.935 & \(+0.26\%\) & 0.117 \\
random\_mask & \(-0.01\%\) & 0.950 & \(-0.29\%\) & 0.090 \\
half\_lambda & \(-0.05\%\) & 0.901 & \(-0.25\%\) & 0.608 \\
no\_wd & \(-0.05\%\) & 0.825 & \(-0.49\%\) & 0.312 \\
\end{longtable}

\textbf{Equivalence testing (TOST).} We apply Two One-Sided Tests (TOST)
with practical equivalence margins \(\delta = \pm 0.5\%\) and
\(\delta = \pm 1.0\%\). At \(\delta = \pm 0.5\%\), only cosine\_schedule
on CIFAR-10 achieves confirmed equivalence
(\(p_{\mathrm{TOST}} = 0.039\)), consistent with limited power at
\(N = 3\). At \(\delta = \pm 1.0\%\), 6 of 12 method--dataset
comparisons confirm equivalence (\(p_{\mathrm{TOST}} < 0.05\)),
including cwd\_hard, cosine\_schedule, random\_mask, and no\_wd on
CIFAR-10. The remaining comparisons are inconclusive.

\textbf{Key interpretation.} The TOST results are intentionally
underpowered at \(N = 3\) seeds, and we present them transparently. The
primary evidence for the Phi Invariance Conjecture comes from the SGD
negative control (Section 6.1): the same \(N = 3\) design that yields
inconclusive TOST results for AdamW yields highly significant
differences (\(p < 0.001\)) for SGD no\_wd. This asymmetry demonstrates
that the AdamW null results reflect a genuine property of the optimizer
rather than universal noise at CIFAR scale.

\hypertarget{budget-equivalence-analysis}{%
\subsubsection{5.2 Budget Equivalence
Analysis}\label{budget-equivalence-analysis}}

Figure 3 plots mean test accuracy against BEM for all seven methods. BEM
values span the full range from 0.0 (constant) to 1.0 (no\_wd). Despite
this ten-fold variation in effective weight decay budget, accuracy
remains essentially flat: 0.25\% spread on CIFAR-10, 0.76\% on
CIFAR-100.

\textbf{A 10\(\times\) variation in effective weight decay budget (BEM
range \(0.0\)--\(1.0\)) produces less than \(0.5\%\) accuracy variation
under AdamW.} This rules out the hypothesis that dynamic weight decay
improves accuracy by using less total weight decay budget.

\hypertarget{csi-and-ais-diagnostic-analysis}{%
\subsubsection{5.3 CSI and AIS Diagnostic
Analysis}\label{csi-and-ais-diagnostic-analysis}}

\textbf{CSI analysis.} Figure 4 and Tables 4a--4b report Coupling
Stability Index values. On CIFAR-10, CSI ranges from 0.838 (SWD) to
0.964 (no\_wd); on CIFAR-100, the range is narrower (0.854--0.868).
\textbf{CSI does not predict accuracy}: Spearman rank correlation
between CSI and accuracy rank is \(\rho < 0.3\) (\(p > 0.3\)) on both
datasets. CSI captures differences in training dynamics but these do not
translate to performance differences under AdamW.

\textbf{AIS analysis.} All methods show AIS values in the range
0.280--0.410. \textbf{AIS is consistent across all methods}, including
random\_mask and no\_wd. The gradient-weight alignment signal is an
\emph{intrinsic property of the network and loss landscape}, not
generated or exploited by the weight decay strategy.

This directly challenges the motivation behind CWD: if AIS is the same
for CWD (which conditions on alignment) and random\_mask (which ignores
alignment), the alignment signal provides no additional useful
information for weight decay decisions on these benchmarks.

\textbf{Table 4a: Diagnostic metrics (CIFAR-10).}

\begin{longtable}[]{@{}lcccc@{}}
\toprule\noalign{}
Method & CSI & AIS & BEM & Weight Norm \\
\midrule\noalign{}
\endhead
\bottomrule\noalign{}
\endlastfoot
constant & 0.841 & 0.336 & 0.000 & 95.89 \\
cosine\_schedule & 0.936 & 0.352 & 0.503 & 96.28 \\
random\_mask & 0.892 & 0.359 & 0.500 & 96.38 \\
half\_lambda & 0.853 & 0.410 & 0.500 & 96.34 \\
no\_wd & 0.964 & 0.343 & 1.000 & 97.04 \\
cwd\_hard & 0.851 & 0.368 & 0.503 & 96.46 \\
swd & 0.838 & 0.360 & 0.900 & 96.84 \\
\end{longtable}

\emph{BEM note for half\_lambda:} \texttt{half\textbackslash{}\_lambda}
is implemented as \(\lambda = 2.5 \times 10^{-4}\) with
\(\varphi \equiv 1\), so the raw computational log yields BEM = 0.000
(no deviation from constant \(\varphi\)). The corrected conceptual BEM =
0.500 (half the total decay budget) is reported in the table and used in
Figure 3.

\textbf{Table 4b: Diagnostic metrics (CIFAR-100).}

\begin{longtable}[]{@{}lcccc@{}}
\toprule\noalign{}
Method & CSI & AIS & BEM & Weight Norm \\
\midrule\noalign{}
\endhead
\bottomrule\noalign{}
\endlastfoot
cosine\_schedule & 0.868 & 0.344 & 0.503 & 105.15 \\
constant & 0.864 & 0.329 & 0.000 & 104.72 \\
swd & 0.854 & 0.339 & 0.900 & 105.78 \\
half\_lambda & 0.866 & 0.342 & 0.500 & 105.45 \\
random\_mask & 0.867 & 0.320 & 0.501 & 105.43 \\
cwd\_hard & 0.855 & 0.321 & 0.500 & 105.49 \\
no\_wd & 0.867 & 0.280 & 1.000 & 106.03 \\
\end{longtable}

\hypertarget{weight-norm-analysis}{%
\subsubsection{5.4 Weight Norm Analysis}\label{weight-norm-analysis}}

Figure 5 shows per-layer mean weight norm trajectories over 200 training
epochs (AdamW, CIFAR-10). Despite the ten-fold BEM variation, all
methods converge to similar final weight norms: 95.89 (constant) to
97.04 (no\_wd), a 1.2\% difference. On CIFAR-100: 104.72 to 106.03, a
1.3\% difference.

This convergence provides the mechanistic explanation for the accuracy
equivalence. AdamW's adaptive step size
\(\eta_t / (\sqrt{\hat{v}_t} + \epsilon)\) scales inversely with
gradient magnitude. When weight decay drives a parameter toward zero,
AdamW's adaptive scaling compensates by increasing the effective step
size. The result is an \textbf{implicit weight norm control mechanism}
that subsumes Phi modulation.

\textbf{Order-of-magnitude argument for the absorption mechanism.} The
ratio of the Phi-modulated perturbation to the gradient update is: \[
\frac{|\lambda \cdot (\varphi - 1) \cdot \|\boldsymbol{\theta}\|| \cdot \sqrt{\hat{v}_t}}{\eta_t \cdot \|\mathbf{g}_t\|} \approx \frac{\lambda \cdot \|\boldsymbol{\theta}\|}{\eta_t / \sqrt{\hat{v}_t}} \approx \frac{5 \times 10^{-4} \times 96}{0.01\text{--}0.1} \approx 0.05\text{--}0.5
\] For maximal modulation (\(|\varphi - 1| \leq 1\)), Phi contributes at
most 5--50\% of the gradient update magnitude, with the higher end
corresponding to late training when \(\eta_t\) is small. At moderate
\(\lambda\) and typical training dynamics, the absorption is effective;
at large \(\lambda\) or very late training, the effect may become
first-order---consistent with the boundary conditions in Section 6.3.

\begin{center}\rule{0.5\linewidth}{0.5pt}\end{center}

\hypertarget{discussion}{%
\subsection{6. Discussion}\label{discussion}}

\hypertarget{the-phi-invariance-conjecture-and-sgd-validation}{%
\subsubsection{6.1 The Phi Invariance Conjecture and SGD
Validation}\label{the-phi-invariance-conjecture-and-sgd-validation}}

Our experimental results reveal a striking pattern: under AdamW, no
weight decay strategy produces statistically distinguishable accuracy
from the constant baseline. We formalize this as:

\textbf{Conjecture 1} (Phi Invariance under AdamW). \emph{For a neural
network trained with AdamW to convergence on CIFAR-scale benchmarks with
BatchNorm and moderate weight decay (\(\lambda \lesssim 10^{-3}\)), the
final test accuracy is invariant to the choice of Phi modulator
\(\varphi(t, \boldsymbol{\theta}, \mathbf{g})\), up to the effective
weight decay budget \(\mathbb{E}[\lambda_{\mathrm{eff}}]\).}

Formally, for any two budget-equivalent modulators \(\varphi_1\) and
\(\varphi_2\): \[
|\mathrm{acc}(\varphi_1) - \mathrm{acc}(\varphi_2)| \leq \epsilon_{\mathrm{noise}}
\] where \(\epsilon_{\mathrm{noise}}\) is bounded by training
stochasticity, not by the functional form of \(\varphi\). The scope is
deliberately restricted to match the experimental evidence; Section 6.3
discusses where the conjecture is expected to fail.

\textbf{SGD negative control: validating the conjecture as
AdamW-specific.} Table 5 presents the SGD control results on CIFAR-10
(49 total SGD experiments). Under SGD, weight decay modulation is NOT
invariant:

\begin{itemize}
\tightlist
\item
  \textbf{no\_wd is significantly worse}: \(-0.91\%\) vs.~constant
  (\(p < 0.001\), two-sample \(t\)-test). Under AdamW, the same no\_wd
  achieves \(-0.05\%\) (\(p = 0.825\)).
\item
  \textbf{SWD is significantly worse}: \(-0.51\%\) vs.~constant
  (\(p = 0.013\)). Under AdamW, SWD achieves \(-0.25\%\)
  (\(p = 0.513\)).
\item
  \textbf{half\_lambda is significantly worse}: \(-0.37\%\)
  (\(p = 0.028\)), confirming that weight decay magnitude matters under
  SGD.
\end{itemize}

\textbf{Table 5: SGD control experiments (CIFAR-10).} Same seven methods
and seeds as the AdamW experiments.

\begin{longtable}[]{@{}lccc@{}}
\toprule\noalign{}
Method & SGD Acc. (\%) & \(\Delta\) vs.~SGD constant & \(p\)-value \\
\midrule\noalign{}
\endhead
\bottomrule\noalign{}
\endlastfoot
constant & \textbf{91.22} \(\pm\) 0.06 & --- & --- \\
cosine\_schedule & 91.20 \(\pm\) 0.10 & \(-0.02\%\) & 0.849 \\
random\_mask & 90.77 \(\pm\) 0.37 & \(-0.44\%\) & 0.168 \\
cwd\_hard & 90.87 \(\pm\) 0.35 & \(-0.35\%\) & 0.238 \\
half\_lambda & 90.84 \(\pm\) 0.15 & \(-0.37\%\) & 0.028* \\
swd & 90.71 \(\pm\) 0.16 & \(-0.51\%\) & 0.013* \\
no\_wd & 90.30 \(\pm\) 0.08 & \(-0.91\%\) & \(<0.001\)** \\
\end{longtable}

*\(p < 0.05\); **\(p < 0.001\).

The critical comparison is cross-optimizer: no\_wd achieves 90.08\%
under AdamW (\(\Delta = -0.05\%\)) vs.~90.30\% under SGD
(\(\Delta = -0.91\%\)). The same \(N = 3\) experimental design detects a
large, highly significant effect under SGD that it cannot detect under
AdamW. This is not a power artifact---the design has 80\% power to
detect effects \(\geq 0.7\%\), and the SGD no\_wd effect of \(0.91\%\)
clears this threshold. The AdamW no\_wd effect of \(0.05\%\) does not.
The SGD results confirm that the Phi Invariance is a genuine property of
AdamW's adaptive mechanism.

\textbf{Mechanistic hypothesis.} AdamW's per-parameter adaptive learning
rate \(\eta_t / (\sqrt{\hat{v}_t} + \epsilon)\) provides implicit
equalization of effective regularization strength across parameters. A
Phi modulator that redistributes weight decay operates on a system that
has already ``pre-equalized'' per-parameter dynamics. SGD, lacking this
pre-equalization, allows Phi modulation to be first-order important. The
mechanistic evidence is:

\begin{enumerate}
\def\labelenumi{\arabic{enumi}.}
\tightlist
\item
  \textbf{Weight norm convergence under AdamW}: All seven methods
  converge to similar weight norms (1.2\% range on CIFAR-10).
\item
  \textbf{Weight norm divergence under SGD}: no\_wd under SGD achieves
  final weight norm 104.6 vs.~64.5 for constant (a 62\% difference),
  consistent with the \(-0.91\%\) accuracy drop.
\item
  \textbf{AIS as an intrinsic property}: AIS is consistent across all
  AdamW methods, confirming the alignment signal is a property of the
  loss landscape not exploitable through weight decay modulation.
\item
  \textbf{Budget insensitivity under AdamW}: Even at BEM = 1.0 (zero
  weight decay), AdamW accuracy is equivalent to constant---consistent
  with AdamW's implicit norm control making explicit decay redundant.
\end{enumerate}

\hypertarget{implications-for-weight-decay-research}{%
\subsubsection{6.2 Implications for Weight Decay
Research}\label{implications-for-weight-decay-research}}

\textbf{For AdamW practitioners.} Under the tested conditions, constant
weight decay with a grid-searched \(\lambda\) is optimal. Engineering
effort on dynamic weight decay strategies offers no measurable benefit
on CIFAR-scale AdamW experiments.

\textbf{For weight decay method developers.} New methods should be
primarily evaluated where the conjecture fails: with SGD (where our data
confirms differences exist), at ImageNet or LLM scale, or in severely
overfitting regimes. Evaluating novel methods solely on CIFAR with AdamW
is misleading---differences are masked by the adaptive absorption
mechanism.

\textbf{For benchmark designers.} This explains why many dynamic weight
decay papers report improvements only in specific settings (CWD with
Lion/Muon, SWD with SGD): benefits are optimizer-specific or
scale-dependent. Our SGD experiments confirm this: SWD and no\_wd, which
are indistinguishable from constant under AdamW, show clear degradation
under SGD.

\textbf{The Phi framework as infrastructure.} Even under the Phi
Invariance Conjecture, the framework provides a common mathematical
language, programmatic interface, and diagnostic metrics. This
infrastructure is essential for future work at the boundary conditions.

\hypertarget{boundary-conditions-when-the-conjecture-may-fail}{%
\subsubsection{6.3 Boundary Conditions: When the Conjecture May
Fail}\label{boundary-conditions-when-the-conjecture-may-fail}}

\begin{itemize}
\tightlist
\item
  \textbf{SGD (confirmed by our experiments).} no\_wd: \(-0.91\%\)
  (\(p < 0.001\)); SWD: \(-0.51\%\) (\(p = 0.013\)). Weight decay
  magnitude and modulation matter under non-adaptive optimizers.
\item
  \textbf{Very large weight decay (\(\lambda \gg 5 \times 10^{-4}\)).}
  The order-of-magnitude argument shows the absorption weakens as
  \(\lambda\) increases; at \(\lambda \sim 10^{-2}\), the Phi
  perturbation becomes first-order comparable to the gradient update.
\item
  \textbf{ImageNet and LLM scale.} At ResNet-50 scale, weight decay's
  dynamics-modifying role may be more consequential; at LLM scale,
  timing effects compound over many tokens.
\item
  \textbf{Severe overfitting regimes.} Our experiments operate in
  well-generalized regimes (\$\sim\$9.7\% generalization gap on
  CIFAR-10). Weight decay's regularization role may reassert in heavily
  overfitted settings.
\item
  \textbf{Vision Transformers with LayerNorm.} LayerNorm interacts
  differently with weight decay than BatchNorm and may respond
  differently to Phi modulation.
\end{itemize}

\hypertarget{cosine-schedule-variance-reduction}{%
\subsubsection{6.4 Cosine Schedule Variance
Reduction}\label{cosine-schedule-variance-reduction}}

cosine\_schedule achieves anomalously low variance on CIFAR-10
(\(\sigma = 0.07\%\) vs.~\(\sigma \approx 0.25\)--\(0.32\%\) for all
other AdamW methods)---a four-fold reduction in seed sensitivity. A
possible mechanistic explanation: the gradual reduction of weight decay
toward the end of training creates a more deterministic convergence
trajectory, reducing stochastic variation from the interplay between
random initialization and weight decay. This effect does not replicate
on CIFAR-100 (\(\sigma = 0.42\%\)), suggesting it may be specific to
lower-complexity settings. If it generalizes to larger scales, it could
be a reproducibility benefit orthogonal to the accuracy null result.

\hypertarget{limitations}{%
\subsubsection{6.5 Limitations}\label{limitations}}

\begin{enumerate}
\def\labelenumi{\arabic{enumi}.}
\item
  \textbf{Statistical power.} Three seeds per configuration provide
  limited power (80\% power to detect \(\geq 0.7\%\) at
  \(\alpha = 0.05\)). The SGD negative control provides design-level
  validation not relying on TOST power.
\item
  \textbf{Scale.} All experiments use ResNet-20 on CIFAR-10/100
  (\$\sim\$270K parameters). The conjecture may not hold at ImageNet
  scale or beyond.
\item
  \textbf{Architecture diversity.} Only ResNet-20 with BatchNorm is
  tested. VGG-16-BN and Vision Transformers may exhibit different
  sensitivity.
\item
  \textbf{Fixed hyperparameters.} Identical hyperparameters ensure fair
  comparison but may disadvantage methods with strong hyperparameter
  sensitivity (e.g., CWD with different \(\beta\), SWD with different
  gradient-norm sensitivity scales).
\item
  \textbf{Single base \(\lambda\).} All experiments use
  \(\lambda = 5 \times 10^{-4}\). The absorption mechanism is expected
  to weaken at larger \(\lambda\).
\end{enumerate}

\begin{center}\rule{0.5\linewidth}{0.5pt}\end{center}

\hypertarget{conclusion}{%
\subsection{7. Conclusion}\label{conclusion}}

We introduced the Phi Modulator Framework, the first unified
mathematical abstraction for dynamic weight decay methods in deep
learning. The framework recovers Cautious Weight Decay, Scheduled Weight
Decay, cosine schedules, Weight Norm Control, AlphaDecay, and all their
compositions as special cases of
\(\varphi(t, \boldsymbol{\theta}, \mathbf{g})\) along four axes:
temporal, directional, spatial, and target-norm. Three diagnostic
metrics---BEM, CSI, and AIS---provide the first quantitative tools for
characterizing weight decay behavior beyond final accuracy.

A systematic benchmark of 42 AdamW experiments finds that all dynamic
weight decay variants are statistically equivalent to constant weight
decay (\(p > 0.05\), including the no-weight-decay case). The \emph{Phi
Invariance Conjecture}---scoped to CIFAR-scale BatchNorm networks under
AdamW with moderate \(\lambda\)---formalizes this: AdamW's adaptive
per-parameter scaling subsumes the effect of any Phi modulator, making
the functional form of \(\varphi\) irrelevant at this scale.

The critical validation comes from 49 SGD control experiments: under
SGD, removing weight decay causes \(-0.91\%\) (\(p < 0.001\)) and
gradient-norm scheduling (SWD) degrades by \(-0.51\%\) (\(p = 0.013\)).
The same experimental design that yields null results under AdamW yields
highly significant differences under SGD, confirming that AdamW
invariance is a genuine property of adaptive optimization, not a
statistical artifact.

Practitioners using AdamW can safely rely on constant weight decay at
CIFAR scale. The most important open question is whether the Phi
Invariance Conjecture holds at ImageNet and LLM scale---the
order-of-magnitude absorption argument suggests it weakens as
\(\lambda\) increases and training duration grows. Violations at larger
scale would establish the precise boundary where dynamic weight decay
begins to matter under AdamW.

\begin{center}\rule{0.5\linewidth}{0.5pt}\end{center}

\hypertarget{references}{%
\subsection{References}\label{references}}

Chen, L., et al.~(2026a). Cautious Weight Decay. \emph{ICLR 2026}.

Chen, L., et al.~(2026b). AdamO: Decoupling Radial and Tangential
Dynamics in Weight Decay. \emph{ICLR 2026}.

D'Angelo, F., et al.~(2024). Why Weight Decay Works: A Comprehensive
Analysis. \emph{NeurIPS 2024}.

Ferbach, D., et al.~(2026). ADANA: Logarithmic-Time Schedules for Weight
Decay and Momentum. \emph{arXiv preprint}.

Fernandez-Hernandez, A., et al.~(2025). The Overfitting-Underfitting
Indicator for Weight Decay Quality. \emph{arXiv preprint}.

Galanti, T., et al.~(2022). SGD and Weight Decay Provably Induce a
Low-Rank Bias in Neural Networks. \emph{NeurIPS 2022}.

He, K., et al.~(2016). Deep Residual Learning for Image Recognition.
\emph{CVPR 2016}.

He, Y., et al.~(2025). AlphaDecay: Module-Wise Weight Decay Guided by
Spectral Density Analysis. \emph{arXiv preprint}.

Kobayashi, G., et al.~(2024). Weight Decay on Multiplicative Parameters
is Nuclear Norm Regularization. \emph{ICML 2024}.

Kosson, A., et al.~(2023). Rotational Equilibrium: How Weight Decay
Balances Learning Across Neural Networks. \emph{NeurIPS 2023}.

Krizhevsky, A. (2009). Learning Multiple Layers of Features from Tiny
Images. Technical Report, University of Toronto.

Krogh, A. \& Hertz, J. (1991). A Simple Weight Decay Can Improve
Generalization. \emph{NeurIPS 1991}.

Loshchilov, I. (2023). AdamWN: Weight Norm Control in Decoupled Weight
Decay. \emph{arXiv preprint}.

Loshchilov, I. \& Hutter, F. (2019). Decoupled Weight Decay
Regularization. \emph{ICLR 2019}.

Tian, Y., et al.~(2024). Selective Projection Decay for
Parameter-Efficient Fine-Tuning. \emph{arXiv preprint}.

Wang, A. \& Aitchison, L. (2024). Optimal Weight Decay as an EMA
Timescale. \emph{ICLR 2024}.

Wilson, A. C., et al.~(2017). The Marginal Value of Adaptive Gradient
Methods in Machine Learning. \emph{NeurIPS 2017}.

Xie, Z., et al.~(2023). Scheduled Weight Decay (AdamS). \emph{AAAI
2023}.

Xie, Y. \& Li, Z. (2024). AdamW Implicitly Performs L-infinity Norm
Constrained Optimization. \emph{ICLR 2024}.

\begin{center}\rule{0.5\linewidth}{0.5pt}\end{center}

\hypertarget{figures-and-tables}{%
\subsection{Figures and Tables}\label{figures-and-tables}}

\begin{itemize}
\tightlist
\item
  Figure 1: fig1\_taxonomy.png --- Phi Modulator Framework taxonomy
  showing four modulation axes (temporal, directional, spatial,
  target-norm) and special cases
\item
  Figure 2: fig2\_accuracy\_comparison.png --- Main accuracy comparison
  bar chart for all seven methods on CIFAR-10 and CIFAR-100 under AdamW
\item
  Figure 3: fig3\_bem\_vs\_accuracy.png --- Budget Equivalence Metric
  vs.~test accuracy scatter plot (CIFAR-10 and CIFAR-100), showing flat
  accuracy across full BEM range
\item
  Figure 4: fig4\_diagnostic\_heatmap.png --- Diagnostic metric heatmap
  (CSI, AIS, BEM) across all methods and both datasets
\item
  Figure 5: {[}TODO: weight norm trajectory figure, generate from
  epoch\_metrics.jsonl{]} --- Per-layer weight norm trajectories over
  200 training epochs for all seven AdamW methods (CIFAR-10)
\item
  Table 1: inline --- Method catalog with Phi expressions and modulation
  axes
\item
  Table 2: inline --- Main AdamW accuracy results (mean ± std over 3
  seeds)
\item
  Table 3: inline --- Statistical tests vs.~constant baseline (AdamW
  paired t-tests)
\item
  Table 4a/4b: inline --- Diagnostic metrics (CSI, AIS, BEM, weight
  norm) for CIFAR-10 and CIFAR-100
\item
  Table 5: inline --- SGD negative control results (CIFAR-10, mean ± std
  over 3 seeds)
\end{itemize}

\end{document}
